% =====================================================================
% Capítulo · Despliegue del Sistema (borrador, fase F7)
% Lo incluye \chapter{Despliegue del Sistema}; aquí solo van secciones.
% Fuentes: apps/web/Dockerfile, docker-compose.yml,
%   docker-compose.supabase-local.yml, .env.docker.example, railway.json,
%   docs/DESPLIEGUE.md y docs/tfg/DECISIONES.md (ADR-011, ADR-022).
% Paquetes necesarios: booktabs, tabularx, listings, tikz
%   (bibliotecas: positioning, arrows.meta, fit).
% =====================================================================

Este capítulo describe cómo se empaqueta y se pone en producción PymeKit. El
objetivo de la fase de despliegue era doble: que cualquier persona pudiera
reproducir la plataforma completa en su equipo con un único comando, y que la
misma configuración sirviera para publicar una instancia de demostración
accesible desde Internet. Ambos objetivos se apoyan en una decisión de diseño
previa: el CMS de datos forma parte de la aplicación web y no es un servicio
aparte (ADR-011), de modo que el sistema que hay que desplegar consta de una
sola aplicación más los servicios gestionados de los que depende.

\section{Arquitectura de despliegue}
\label{sec:despliegue-arquitectura}

La Figura~\ref{fig:despliegue} muestra los nodos que intervienen en una
instalación de PymeKit. El único componente propio es el servidor de la
aplicación, un proceso Node.js sin estado que sirve tanto las páginas
renderizadas en el servidor como las rutas de API (\emph{server functions},
\emph{webhooks} y la API del CMS). Todo el estado persistente reside en la
plataforma de base de datos como servicio (Postgres, autenticación y
almacenamiento de ficheros), y los pagos se delegan en la pasarela externa.

\begin{figure}[htbp]
  \centering
  % Se escala al ancho de línea para que encaje en cualquier plantilla.
  \resizebox{\textwidth}{!}{%
  \begin{tikzpicture}[
      node distance=1.4cm and 2.2cm,
      box/.style={draw, rounded corners, align=center, minimum width=3.4cm,
                  minimum height=1.1cm, font=\small},
      group/.style={draw, dashed, rounded corners, inner sep=0.35cm},
      arr/.style={-{Stealth[length=2.2mm]}, thick},
      lbl/.style={font=\scriptsize, midway, fill=white, inner sep=1pt}
    ]
    \node[box] (browser) {Navegador\\(usuario o administrador)};
    \node[box, right=of browser] (web) {Servidor de la aplicación\\Node.js (contenedor)\\\texttt{/}, \texttt{/admin/cms}, \texttt{/api/*}};
    \node[box, above right=0.6cm and 2.2cm of web] (db) {Postgres\\(RLS, esquema \texttt{cms})};
    \node[box, right=2.2cm of web] (auth) {Auth y Storage};
    \node[box, below right=0.6cm and 2.2cm of web] (stripe) {Pasarela de pagos\\(modo test)};
    \node[group, fit=(db)(auth), label={[font=\scriptsize]above:Plataforma BaaS}] {};

    \draw[arr] (browser) -- node[lbl]{HTTPS} (web);
    \draw[arr] (web) -- node[lbl, sloped]{SQL / REST} (db);
    \draw[arr] (web) -- node[lbl]{REST} (auth);
    \draw[arr] (web) -- node[lbl, sloped]{API} (stripe);
    \draw[arr, dashed] (stripe.west) to[bend left=25] node[lbl, sloped]{\emph{webhook}} (web.south);
    \draw[arr, dashed] (browser.north) to[bend left=20] node[lbl, sloped]{sesión (clave pública)} (auth.north west);
  \end{tikzpicture}%
  }
  \caption{Diagrama de despliegue de PymeKit.}
  \label{fig:despliegue}
\end{figure}

La Tabla~\ref{tab:componentes-despliegue} resume la función de cada componente
y el artefacto que lo configura dentro del repositorio.

\begin{table}[htbp]
  \centering
  \small
  \begin{tabularx}{\textwidth}{@{}lXl@{}}
    \toprule
    \textbf{Componente} & \textbf{Función} & \textbf{Configuración} \\
    \midrule
    Servidor de la aplicación & Renderizado, \emph{server functions}, API del CMS y \emph{webhooks} & \texttt{apps/web/Dockerfile} \\
    Orquestación local & Construye la imagen y la arranca con sus variables & \texttt{docker-compose.yml} \\
    Base de datos & Esquemas, políticas RLS y RBAC del CMS & \texttt{apps/web/supabase/migrations} \\
    Autenticación y ficheros & Sesiones, MFA, URLs de redirección, \emph{buckets} & \texttt{config.toml} / panel \\
    Pasarela de pagos & Suscripciones; notifica cambios por \emph{webhook} & variables \texttt{STRIPE\_*} \\
    Plataforma de alojamiento & Build, arranque y sonda de salud en la nube & \texttt{railway.json} \\
    \bottomrule
  \end{tabularx}
  \caption{Componentes del despliegue y su configuración.}
  \label{tab:componentes-despliegue}
\end{table}

\section{Contenerización de la aplicación}
\label{sec:despliegue-imagen}

La aplicación se empaqueta en una imagen de contenedor definida mediante un
\emph{Dockerfile} de tres etapas. Construir por etapas permite que las
herramientas de compilación, que ocupan cientos de megabytes, no lleguen a la
imagen final:

\begin{enumerate}
  \item \textbf{Recorte del monorepo.} La herramienta de orquestación del
    monorepo calcula el subconjunto de paquetes del que depende la aplicación
    web y separa sus manifiestos del resto del código. Así, la capa que
    instala dependencias solo se invalida cuando cambian estas, y no con cada
    modificación del código fuente.
  \item \textbf{Compilación.} Se instalan las dependencias a partir del
    fichero de bloqueo (instalación reproducible) y se ejecuta la build de
    producción. El servidor resultante queda autocontenido: el empaquetador
    incluye en la carpeta de salida las librerías que el servidor necesita,
    de forma que no hace falta copiar el directorio de dependencias.
  \item \textbf{Ejecución.} La imagen final parte de una distribución mínima
    de Node.js y solo contiene la carpeta de salida. El proceso se ejecuta
    con un usuario sin privilegios y la imagen declara una sonda de salud
    sobre la ruta \texttt{/api/healthcheck}, que el motor de contenedores usa
    para marcar el servicio como sano o reiniciarlo.
\end{enumerate}

Un aspecto que condiciona el diseño es la distinción entre variables
\emph{de compilación} y \emph{de ejecución}. Las variables públicas (prefijo
\texttt{VITE\_}: URL del sitio, URL y clave pública de la base de datos, clave
publicable de la pasarela) se incrustan en el código durante la build, por lo
que se pasan como argumentos de construcción. Los secretos de servidor (clave
secreta de la base de datos, cadena de conexión de la API del CMS, claves de
la pasarela y del correo) se inyectan únicamente al arrancar el contenedor y
nunca quedan dentro de la imagen, lo que permite publicarla o reutilizarla sin
exponer credenciales. Para que los argumentos vacíos no anulen los valores por
defecto versionados en el repositorio, la build solo escribe los que llegan
con valor y falla de forma explícita si falta alguno de los obligatorios.

\section{Orquestación con Docker Compose}
\label{sec:despliegue-compose}

El fichero \texttt{docker-compose.yml} de la raíz describe el servicio de la
aplicación. No levanta la base de datos: la aplicación se conecta a una
instancia ya existente, sea la del entorno de desarrollo local o un proyecto
en la nube, y el mismo fichero sirve para ambos casos. El
Listado~\ref{lst:compose} recoge su parte esencial.

\begin{lstlisting}[language=bash, caption={Servicio de la aplicación en \texttt{docker-compose.yml} (extracto).}, label={lst:compose}, basicstyle=\ttfamily\footnotesize]
services:
  web:
    image: pymekit-web:latest
    build:
      context: .
      dockerfile: apps/web/Dockerfile
      args:            # publicas: se incrustan en la build
        VITE_SITE_URL: ${VITE_SITE_URL:?Falta VITE_SITE_URL}
        VITE_SUPABASE_URL: ${VITE_SUPABASE_URL:?Falta VITE_SUPABASE_URL}
        VITE_SUPABASE_PUBLIC_KEY: ${...:?Falta}
    env_file:          # de servidor: se leen al arrancar
      - path: .env.docker
        required: true
    ports:
      - '${WEB_PORT:-3200}:3000'
    restart: unless-stopped
\end{lstlisting}

Todas las variables se documentan en una plantilla versionada,
\texttt{.env.docker.example}, sin valores secretos; el fichero real con las
credenciales está excluido tanto del control de versiones como del contexto
de construcción de la imagen. Para el entorno local se añade un segundo
fichero de ajuste que ejecuta el contenedor con la red del anfitrión. La razón
es que la URL de la base de datos se incrusta en la build y la usan tanto el
navegador como el servidor: con la red del anfitrión, la dirección local
designa el mismo servicio para ambos y no es necesario mantener dos URLs
distintas.

\section{Despliegue en la nube}
\label{sec:despliegue-nube}

La instancia de demostración combina tres servicios gestionados:

\begin{itemize}
  \item \textbf{Base de datos como servicio.} Se crea un proyecto y se le
    aplican las migraciones del repositorio con la interfaz de línea de
    órdenes de la plataforma (\texttt{supabase db push}), que incluyen tanto
    el esquema de la aplicación como el del CMS. Después se configuran las
    URLs autorizadas de autenticación, un servidor de correo propio y el
    \emph{webhook} de base de datos que avisa a la aplicación cuando se borra
    una suscripción. Los datos de ejemplo (\emph{seed}) no se aplican en este
    entorno.
  \item \textbf{Pasarela de pagos en modo test.} Se registra un
    \emph{webhook} público que apunta a \texttt{/api/billing/webhook}; la
    aplicación verifica la firma de cada evento con el secreto compartido
    antes de actualizar las suscripciones en la base de datos.
  \item \textbf{Plataforma de alojamiento.} El fichero \texttt{railway.json}
    declara la orden de build, la de arranque, la ruta de la sonda de salud y
    la política de reinicio ante fallos. Las variables se definen en el panel
    del servicio y cada integración en la rama principal produce un nuevo
    despliegue.
\end{itemize}

El procedimiento completo, junto con la actualización y las copias de
seguridad, se recoge en el manual de instalación (Anexo~\ref{anx:manual-instalacion}).
% PENDIENTE: crear la etiqueta anx:manual-instalacion al pasar
% docs/DESPLIEGUE.md al anexo de la memoria (F8).

\section{Verificación}
\label{sec:despliegue-verificacion}

La imagen se construyó y se arrancó contra la base de datos del entorno local,
comprobando que la página de inicio y la de inicio de sesión responden con
código 200 y que la sonda de salud informa de que la base de datos es
accesible. La imagen final ocupa unos 265~MB en disco
(medido el 20/10/2026), frente a los varios gigabytes que suponen las dependencias de
desarrollo del monorepo.

% PENDIENTE: completar con la URL y la fecha de la instancia de demostración
% cuando se publique (F7, tareas 3–5), y con la prueba del manual desde cero.
